\section{Utilizzo del codice}
\label{sec:utilizzo}

\subsection{Contenuto della consegna}

\begin{lstlisting}
progetto_ASD/
|-- README.md
|-- pom.xml                   build Maven
|-- asd-project.jar           eseguibile precompilato
|-- requirements-report.txt   dipendenze Python per i grafici
|-- src/main/java/mio/        codice applicativo
|-- src/test/java/mio/        test JUnit
|-- risultati/
|   |-- esperimenti.csv       piano sperimentale (67 istanze, 215 righe)
|   |-- riassunti.txt         riassunto di ogni riga del piano
|   |-- limiti_scala.csv      serie sui limiti di scala (108 coppie per variante)
|   |-- griglie_json/         23 griglie del piano
|   |-- grafici/              grafici della relazione
|   `-- genera_grafici.py     script di analisi
`-- relazione/                sorgenti LaTeX e PDF di questa relazione
\end{lstlisting}

\subsection{Requisiti}

\begin{itemize}
  \item Java 17 o successivo per eseguire il JAR;
  \item Maven 3.9 o successivo per compilare ed eseguire i test;
  \item Python 3 con i pacchetti di \code{requirements-report.txt} per
        rigenerare i grafici;
  \item una distribuzione \LaTeX{} con \code{latexmk} per ricompilare la
        relazione.
\end{itemize}

\subsection{Esecuzione}

Tutti i comandi vanno eseguiti dalla radice della consegna: i file di
output sono scritti nella cartella \code{risultati/}.

\paragraph{Piano sperimentale.} Senza argomenti l'applicazione esegue il
piano della \cref{sec:protocollo}, stampa un riassunto per istanza e
variante e riscrive \code{esperimenti.csv}, \code{riassunti.txt} e le griglie
della cartella \code{griglie\_json/}. Richiede meno di un minuto.

\begin{lstlisting}
java -jar asd-project.jar
\end{lstlisting}

\paragraph{Serie sui limiti di scala.} Riscrive
\code{risultati/limiti\_scala.csv} per le varianti \code{base} e
\code{limite-propagato}. Richiede circa 18 minuti, quasi tutti dovuti alle
esecuzioni che raggiungono il limite di 10 s; indicando una variante si
esegue la serie solo per quella.

\begin{lstlisting}
java -jar asd-project.jar --limiti-scala
java -jar asd-project.jar --limiti-scala limite-propagato
\end{lstlisting}

\paragraph{Compito Due.} Calcola contesto, complemento e frontiera di $O$ e
la distanza libera fra $O$ e $D$ su una griglia salvata
(\cref{sec:applicazione-chiusura}).

\begin{lstlisting}
java -jar asd-project.jar --grid-input risultati/griglie_json/scenario_02_12x12_15ost_seed50.json --origin 0,0 --destination 10,10 --chiusura
\end{lstlisting}

\paragraph{Compito Tre, istanza singola.} Genera una griglia con i
parametri indicati, oppure la legge da un file JSON, esegue \cmin{} nei due
versi una sola volta con la variante scelta e salva il risultato. Se $O$ o
$D$ cadono su un ostacolo viene usata la cella libera più vicina e
l'applicazione lo segnala.

\begin{lstlisting}
java -jar asd-project.jar --width 20 --height 20 --obstacles 40 --seed 77 --origin 0,0 --destination 19,19 --summary true --output risultati/prova.csv --grid-output risultati/prova.json
java -jar asd-project.jar --grid-input risultati/prova.json --origin 0,0 --destination 19,19 --variante limite-propagato
\end{lstlisting}

\begin{table}[H]
\centering
\small
\begin{tabularx}{\textwidth}{@{}l X@{}}
\toprule
\textbf{Opzione} & \textbf{Significato e valore predefinito} \\
\midrule
\code{-{}-width N}, \code{-{}-height N} & dimensioni della griglia generata (20 e 20) \\
\code{-{}-obstacles N}   & numero di ostacoli, da 0 a $WH-1$ (20) \\
\code{-{}-seed N}        & seed del generatore (42) \\
\code{-{}-grid-input f.json} & griglia letta da file invece che generata;
                           esclude le quattro opzioni precedenti e
                           \code{-{}-grid-output}; nel CSV il seed vale $-1$ \\
\code{-{}-origin x,y}    & origine $O$ (\code{0,0}) \\
\code{-{}-destination x,y} & destinazione $D$ (angolo opposto all'origine) \\
\code{-{}-chiusura}      & esegue solo il Compito Due: mappa della chiusura di $O$,
                           conteggi e $\dlib(O,D)$ \\
\code{-{}-chiusura-output f.json} & con \code{-{}-chiusura}, salva gli stati delle celle \\
\code{-{}-variante v}    & variante di \cmin{}: \code{base} (predefinita),
                           \code{senza-cache}, \code{senza-ordinamento},
                           \code{limite-propagato} \\
\code{-{}-summary true|false} & stampa del riassunto (\code{true}) \\
\code{-{}-output f.csv}  & CSV del risultato (\code{risultati/esperimento\_personalizzato.csv}) \\
\code{-{}-grid-output f.json} & griglia generata (\code{risultati/griglie\_json/personalizzata.json}) \\
\code{-{}-limiti-scala [v]} & serie sui limiti di scala, per entrambe le
                           varianti principali o solo per $v$; va usata da sola \\
\code{-{}-help}          & elenco delle opzioni \\
\bottomrule
\end{tabularx}
\caption{Opzioni della riga di comando.}
\label{tab:opzioni}
\end{table}

\subsection{Compilazione e test}

\begin{lstlisting}
mvn -q clean package
java -jar target/asd-project-1.0-SNAPSHOT.jar
\end{lstlisting}

Il primo comando compila il codice, esegue i 68 test e crea in
\code{target/} un JAR equivalente a \code{asd-project.jar}. Per eseguire
solo i test si usa \code{mvn -q test}.

\subsection{Grafici e relazione}

\begin{lstlisting}
python -m pip install -r requirements-report.txt
python risultati/genera_grafici.py
cd relazione
latexmk -pdf Relazione_ASD_2024_25.tex
\end{lstlisting}

Lo script dei grafici verifica la struttura dei due file CSV e la coerenza
delle lunghezze fra varianti prima di disegnare, scrive le immagini in
\code{risultati/grafici/}, da cui la relazione le include direttamente, e
stampa i valori riassuntivi riportati nella \cref{sec:compito4}.
